\chapter{Discussion and Limitations}
\label{ch:joint-discussion}

\section{Common Methodological Principle}
\label{sec:joint-shared}

The two prediction tasks investigated in this thesis address different objects, but are
based on a common methodological principle: structured graph knowledge and explicit
empirical statistical evidence provide different, potentially complementary descriptions
of the same system.

In Task~A, the prediction object is a candidate directed relation. The structural branch
describes the position of the candidate nodes within the heterogeneous graph, while the
statistical branch describes empirical dependence patterns observed in the patient
population. In Task~B, the graph topology is shared across patients and the prediction
object becomes a patient-specific clinical endpoint. Patient observations are propagated
through the graph, while an explicit statistical evidence branch provides an additional
description of endpoint-related dependencies.

The common modelling template can therefore be summarized as
\begin{equation}
\text{structured graph representation}
\;+\;
\text{explicit empirical evidence}
\;\rightarrow\;
\text{task-specific prediction}.
\end{equation}
The relative contribution of the two components is task dependent. Statistical evidence is
particularly useful when relevant information is not adequately represented or preserved
by the graph pathway, whereas a sufficiently effective graph encoder may already capture
part of the same signal.
On the synthetic Task~B benchmark this distinction is concrete: when residual connections
preserve endpoint-relevant information inside the message-passing stack, the evaluated
statistical-evidence variants do not provide a consistent additional improvement over the
corresponding graph-only residual models, whereas the same branch recovers substantial
performance when deep models degrade without residual propagation.

\section{What the Two Synthetic Tasks Show}
\label{sec:joint-synthetic}

\subsection*{Task~A: confronting structured knowledge with observational evidence}

Task~A demonstrates that patient-derived statistical evidence can provide substantial
information beyond heterogeneous graph topology alone. Progressive enrichment of the
candidate-specific statistical representation improves directed edge ranking relative to the
graph-only HGT baseline.

Importantly, statistical dependence does not determine edge direction. Direction is supplied
by the ordered candidate definition and the mechanistic graph schema. HCR-derived
coefficients provide empirical evidence about relationships between variables, but are not
interpreted as causal effects.
Task~A can therefore be understood as confronting prior structured knowledge with
population-level observational evidence under a controlled reconstruction objective.
It does not perform causal discovery.

Residual connections inside the Task~A HGT encoder are standard skip connections with
LayerNorm; they are architectural details of the structural branch and are not studied as an
anti-oversmoothing intervention in that chapter.

\subsection*{Task~B: structural context for patient-level endpoints}

Task~B examines whether a shared mechanistic graph can provide useful structural context
for predicting patient-level clinical endpoints.

The experiments show that the usefulness of the statistical branch depends strongly on the
state of the graph propagation pathway. Without residual connections, deeper models can
lose endpoint-relevant information, while the explicit statistical branch can provide a
direct alternative information route. When residual propagation preserves the graph signal,
the statistical branch does not provide a consistent additional improvement on the
synthetic benchmark.
With residual connections enabled, the evaluated statistical-evidence variants did not
provide a consistent additional improvement over the corresponding graph-only residual
models on the synthetic benchmark.

\section{Cross-Domain Proxy Experiments}
\label{sec:joint-proxies}

The proxy experiments test whether the architectural principle survives substantial
changes in graph semantics and prediction objective rather than establishing direct domain
transfer.

Last-FM* replaces pharmacotherapy variables with users, items and knowledge-graph
entities and changes the objective from directed graph reconstruction to recommendation
ranking. The development ablation shows that candidate-specific empirical context
substantially improves the neural HGT model.
The matched sealed benchmark additionally
shows that improving a neural graph model does not imply superiority over strong
classical recommenders:
development gains over HGT are not equivalent to sealed-test dominance, and the proposed
full model remains below the strongest classical baseline (RP$3\beta$) under the sealed
external protocol.

The three Last-FM* result levels therefore answer different questions and must not be
interchanged.
Sampled validation NDCG@20 is a cheap checkpoint-selection proxy;
full-catalog development ranking isolates the contribution of candidate-specific context
relative to HGT alone;
only the sealed external evaluation is comparable to published recommender numbers under
the same evaluator.
In particular, the LEG residual in Last-FM* is an additive correction in decoder score space
and must not be conflated with residual connections inside a message-passing stack:
on the sampled checkpoint metric LEG changes NDCG@20 by only about $0.0004$, whereas on
full-catalog development ranking it raises mean NDCG@20 from about $0.232$ (HGT+A5+H3)
to about $0.276$ (full model).

The Hetionet-derived proxy performs the analogous stress test for Task~B. It preserves the
endpoint-prediction logic while introducing different graph semantics and longer,
partially latent information pathways. Under this setting, explicit statistical evidence
becomes particularly useful for deeper models, and statistical evidence and residual
propagation can become complementary:
at six layers, test AUROC rises from $0.585$ for the graph-only model to $0.757$ with
statistical evidence alone and $0.760$ when both mechanisms are combined.
Hetionet remains a methodological proxy rather than clinical validation.

\section{Prediction versus Mechanistic Clinical Reasoning}
\label{sec:joint-mechanistic-reasoning}

Predictive performance alone does not fully determine the clinical usefulness of a model.
A patient-level predictor primarily addresses the question of whether a clinical event is
likely to occur. Clinical decision making additionally requires understanding which observed
factors and plausible mechanisms may contribute to that risk and which of them could
potentially be modified.

This distinction is especially important in pharmacotherapy. The same clinical endpoint may
arise through different combinations of drug exposure, comorbidity, physiological state,
laboratory abnormalities and intermediate adverse effects. A predicted risk of an endpoint
is therefore potentially more informative when it can be related to plausible mechanistic
routes than when it is presented as an isolated score.

The graph-based formulation investigated in this thesis provides a structural basis for such
an extension. The present work evaluates graph reconstruction and endpoint prediction;
it does not yet identify patient-specific causal mechanisms. A natural next stage would be
to determine which regions and pathways of the mechanistic graph are supported by the
observations of an individual patient.

Instead of treating the predicted endpoint as the final
output, prediction becomes one component of a broader reasoning process:
\begin{equation}
\text{endpoint risk}
\;\rightarrow\;
\text{plausible mechanistic routes}
\;\rightarrow\;
\text{modifiable factors}
\;\rightarrow\;
\text{possible intervention hypotheses}.
\end{equation}
The present architecture does not perform intervention analysis; the sequence above is a
motivating interpretation for future work.

\section{Observational Data, Explainability and Causal Interpretation}
\label{sec:joint-observational-causal}

Clinical data are generated
through a healthcare process that includes decisions about testing, treatment,
documentation and follow-up. Consequently, recorded associations may reflect both the
patient's underlying condition and institution-specific patterns of observation or clinical
practice.

This issue motivates the controlled observational scenarios used in the synthetic benchmark,
including selection bias, noisy documentation and multi-hospital variation. In the synthetic
setting these processes can be modified while the underlying graph remains fixed. In real
clinical data, the corresponding observation mechanisms would be substantially more
complex and only partially observable.

Mechanistic prior knowledge may provide useful structure when empirical associations are
unstable across institutions or populations. It does not remove confounding or establish
causal direction, but it provides an explicit hypothesis about how variables are related that
can be evaluated against observations.

Model explainability and mechanistic explanation should also remain conceptually distinct.
Feature attribution can indicate which model inputs contributed to a prediction, but this
does not establish that those variables constitute the mechanism responsible for the
clinical event. A mechanistic graph provides a different form of explanation by explicitly
representing hypothesized relationships and pathways; causal interpretation still requires
additional assumptions and validation.

The present thesis therefore remains at the level of predictive modelling supported by
mechanistic priors and empirical dependence. Future work could connect this framework
with dedicated causal-discovery or causal-inference methods, but such extensions should
not be inferred from predictive performance alone.

\section{Average Precision and Asymmetric Error Costs in Graph Reconstruction}
\label{sec:joint-ap-reconstruction}

Average Precision provides a threshold-independent summary of the trade-off between
recovering existing relations and avoiding unsupported ones.
It is well suited to model comparison and checkpoint selection in the present
link-prediction setting.
However, it does not explicitly assign different downstream costs to false-positive and
false-negative errors.

This distinction becomes important when link prediction is interpreted as graph
reconstruction rather than only as binary classification.
A false-negative prediction leaves a potentially valid relation unresolved.
A false-positive prediction introduces into the reconstructed graph a relation that is absent
from the audited reference structure.

This difference can be expressed conceptually as
\begin{equation}
\label{eq:joint-cfp-cfn}
C_{\mathrm{FP}}>C_{\mathrm{FN}},
\end{equation}
when the predicted graph is intended to support subsequent structural or mechanistic
reasoning.
The inequality is not imposed by the current training objective; rather, it represents a
task-dependent consideration for downstream use of reconstructed edges.

A missed edge preserves uncertainty about a part of the graph.
A spurious edge can instead convert uncertainty into an incorrect structural assertion and
may alter the set of paths and neighbourhoods available for subsequent reasoning.
In applications in which predicted relations are incorporated into a knowledge graph, it may
therefore be preferable to recover a smaller set of high-confidence edges rather than to
maximize recall at the cost of introducing unsupported structure.

AP evaluates the overall quality of the candidate-edge ranking, but it does not determine the
operating point at which a predicted edge should be accepted into an expanded graph.
Future graph-completion studies could therefore complement AP with precision-oriented
operating-point measures, for example precision at a prescribed recall level, precision among
the highest-ranked candidates, or a validation-selected threshold $\tau$ in a selective
acceptance rule
\begin{equation}
\label{eq:joint-eadd}
\hat{E}_{\mathrm{add}}
=
\bigl\{(A,G):s_{AG}\ge\tau\bigr\}.
\end{equation}
Consequently, graph completion intended for knowledge enrichment should ultimately be
treated not only as a ranking problem, but also as a selective knowledge-acceptance
problem. A high-scoring candidate should become an asserted graph relation only when the
available evidence satisfies an acceptance criterion appropriate to its downstream use.
The present study does not implement such a deployment rule.

\section{Limitations and Deployment Implications}
\label{sec:joint-limitations}

The primary limitation of the present study is the synthetic nature of the pharmacotherapy
benchmark. The known graph and data-generating process provide a controlled environment
for methodological evaluation, but cannot reproduce the full heterogeneity, measurement
noise and confounding of real clinical records.
At the same time, the controlled setting is useful for studying architectural questions that
would be difficult to isolate when both the graph structure and observation process are
unknown. The synthetic benchmark should therefore be interpreted as methodological
ground truth rather than as a substitute for clinical evidence.

The Last-FM* and Hetionet experiments extend the methodological stress testing beyond
the original synthetic setting, but neither constitutes validation on real patient data.
Generator-based oracle references in Task~B are synthetic artefacts and should not be read
as clinical ceilings or as strict theoretical upper bounds.

A practical feature of the proposed framework is its modularity. The graph encoder and
explicit statistical branch can be modified independently, allowing alternative
message-passing operators, dependence estimators or domain-specific evidence channels
to be evaluated without redesigning the complete architecture.

The explicit statistical evidence branch is comparatively lightweight once cohort-level
statistics have been precomputed, although the computational cost of the complete system
depends on graph size, message-passing architecture and the number of candidate
predictions required.

Separating structural and statistical evidence also improves inspectability. Individual
dependence coefficients, evidence blocks and learned gates can be examined independently
of the latent GNN representation. This does not make the complete model intrinsically
interpretable, but provides an explicit evidence layer that can support auditing and
hypothesis generation.

The patient-specific graph formulation can also be viewed as one possible methodological
building block toward future mechanistic digital-twin systems. The present architecture is
not a digital twin: it lacks temporal state updating, validated intervention models,
counterfactual estimation and real-world clinical calibration. Nevertheless, representing
patient observations on a shared mechanistic graph provides a possible intermediate step
between static risk prediction and dynamic patient-specific mechanistic modelling.
